\documentclass[11pt,a4paper]{article}
\usepackage[UTF8]{ctex}
\usepackage{geometry}
\geometry{left=2.5cm,right=2.5cm,top=2.5cm,bottom=2.5cm}
\usepackage{amsmath,amssymb,amsthm}
\usepackage{hyperref}
\usepackage{booktabs}
\usepackage{enumitem}
\usepackage{xltabular}
\hypersetup{colorlinks=true,linkcolor=blue,urlcolor=blue}

\newtheorem{definition}{定义}[section]
\newtheorem{proposition}{命题}[section]
\newtheorem{remark}{注记}[section]

\title{UEBA与对手行为分析的桥接（修正版）\\
——基线粒度的最优裁定、行为基线的柯克霍夫保护与自噬效应的规避}
\author{李龙胜}
\date{\today}

\begin{document}

\maketitle

\begin{center}
\textit{
对手行为分析不再靠模糊的统计信号。\\
UEBA将感知落地为用户和实体的行为基线，检测每一次基线偏离。\\
但基线太粗抓不到攻击，太细则误报淹没分析能力。\\
攻击方不仅跨周期积累知识，还能在同一周期内自适应校准行为。\\
最优基线粒度恰好是边际检测收益等于边际误报成本的均衡点。\\
基线参数需要柯克霍夫保护——周期性扰动对抗跨周期探测，随机噪声对抗单周期试探。\\
防护层过强会吞噬系统分析能力——这是自噬效应，够用就行是唯一解药。\\
本文完成UEBA与四层分析能力分配框架的完整对接。
}
\end{center}

\vspace{5mm}

\begin{abstract}
在已有的四层分析能力分配框架中，
对手行为分析层通过 \(\Delta D, \Delta N, \Delta \tau\) 三个统计信号感知攻击方探测强度。
UEBA的引入将这一抽象层降落到用户和实体的具体行为基线。
本文建立UEBA基线粒度的成本模型——基线越细，检测越灵敏但维护成本和误报成本越高。
推导最优基线粒度的边际均衡条件，并讨论U型安全损失函数成立的边界条件。
分析UEBA对攻击方生态识别、逼移时机感知和蜜罐引流时机的增益。
将UEBA基线参数纳入柯克霍夫原则的保护范围，
周期性扰动对抗跨周期知识积累，随机噪声对抗单周期内自适应校准。
处理UEBA误报的负反馈循环——用贪心择优裁定偏离度筛选的最优阈值。
识别防护层自噬效应的普遍规律，指出够用就行原则是克制自噬的唯一手段。
完成UEBA与四层分配框架的完整对接。
\end{abstract}

\tableofcontents

\section{引言}

在已有的四层分析能力分配框架中，
对手行为分析通过统计信号 \(\Delta D, \Delta N, \Delta \tau\) 感知攻击方探测强度。
这三个信号有效但粒度粗——它们告诉你“探测强度在上升”，
但不能告诉你“哪个内部用户的行为模式发生了变化”。

UEBA填补的正是这个缺口。
UEBA为每个用户和实体建立行为基线——
用户基线的维度包括登录时间、登录地点、常用命令、
文件访问模式、数据外传频率。
实体基线的维度包括服务器的网络连接模式、
进程启动序列、端口监听变化、资源使用周期。

当攻击方使用被盗凭据登录、执行异常命令、访问非惯常文件时，
UEBA检测到的不是统计信号的波动，而是具体行为基线的偏离。
这正是对手行为分析层从抽象概念落地为具体技术的关键一步。

\section{最优基线粒度}

UEBA面临一个核心的资源分配问题——基线粒度。
基线太粗，检测不灵敏，攻击方在基线容忍范围内活动而不触发告警。
基线太细，计算和存储成本急剧上升，误报率同步上升，
分析师的研判精力被系统自身的噪声消耗。
最优基线粒度是够用就行原则在这一维度上的直接应用。

\subsection{基线粒度的定义}

\begin{definition}[基线粒度]
基线粒度 \(g \in [0,1]\) 度量UEBA对用户和实体行为的建模精细度。
\(g = 0\)：极粗基线——只记录用户的登录城市和登录时间。
\(g = 1\)：极细基线——记录每次击键间隔、鼠标移动轨迹、
进程间调用的时序模式。
\end{definition}

\subsection{成本函数}

基线粒度的维护成本随粒度提升而递增且边际成本递增。
设
\[
C_{\text{baseline}}(g) = \kappa_{\text{ueba}} \cdot g^2,
\]
其中 \(\kappa_{\text{ueba}}\) 是UEBA基线运维成本系数。

细粒度基线需要更高频的数据采集、更复杂的模型训练和更大的存储空间，
计算成本超线性增长。二次函数反映这一边际成本递增特性。

基线粒度影响安全损失。极粗基线几乎检测不到异常，漏报成本极高。
极细基线过于灵敏，正常用户的合法波动也触发告警，误报成本上升。
因此预期安全损失函数 \(L(g)\) 在粗粒度端随 \(g\) 增大而下降（检测能力提升），
在细粒度端随 \(g\) 增大而上升（误报淹没分析能力）。

\subsection{最优基线粒度}

总成本为基线维护成本与安全损失之和：
\[
J(g) = \kappa_{\text{ueba}} \cdot g^2 + L(g).
\]
对 \(g\) 求导并令其为零，得到最优基线粒度的均衡条件：
\[
2\kappa_{\text{ueba}} \cdot g^* = -\frac{dL(g)}{dg}\bigg|_{g=g^*}.
\]

最优基线粒度的含义很直观——
不是追求最细的基线，而是追求刚好能捕捉关键攻击信号、
同时不让误报消耗过多分析能力的基线。
多余粒度浪费算力，不够粒度制造盲区。
这是够用就行原则在行为基线维度上的精确表达。

\textbf{U型安全损失函数成立的边界条件}：
安全损失函数的U型假设在漏报损失和误报损失具有可比量级时成立。
当漏报损失远大于误报损失时，
U型曲线的右半段可能永远不会上升——
最优基线粒度可能趋向极细端，
此时计算成本和存储成本（而非误报成本）成为限制因素。
该边界条件在当前模型中已识别，详细定量分析待后续展开。

\section{UEBA与四层分配框架的对接}

引入UEBA后，对手行为分析层的算力消耗被显著细化。

基线粒度的调整消耗算力，直接计入对手行为分析 \(r_{\text{opponent}}\) 的预算。
基线偏离事件的评估需要分析师研判，消耗告警分析 \(r_{\text{alert}}\) 的分析能力。
两个层面通过基线粒度产生耦合——
基线粒度越细，产生的基线偏离事件越多，
需要调配到告警分析层的分析能力也随之增加。

最优基线粒度本身需要被更新，消耗参数重标定 \(r_{\text{param}}\) 的分析能力。
当攻击方探测强度上升时，UEBA基线粒度应自动缩短，
此时安全参数调节层 \(r_{\text{sec}}\) 将分析能力调入基线模型更新，
临时提升基线粒度。

四层之间的最优分配由贪心择优统一裁定：
\[
\frac{\partial U_{\text{alert}}}{\partial r_{\text{alert}}} =
\frac{\partial U_{\text{param}}}{\partial r_{\text{param}}} =
\frac{\partial U_{\text{opponent}}}{\partial r_{\text{opponent}}} =
\frac{\partial U_{\text{sec}}}{\partial r_{\text{sec}}}.
\]

\textbf{技术注释：UEBA运维成本的多维拆分}。
UEBA基线运维成本在工程实现中包含多个子项——
模型训练消耗GPU/CPU算力，与告警分析共享算力预算；
长期行为数据存储，与台账和历史日志共享存储预算；
数据采集消耗网络带宽，与SIEM日志采集共享管道预算。
各子项的成本结构和边际递增特性可能不同。
当前版本以单一系数汇总，在原型部署中应根据实际资源消耗拆分为独立子项并分别标定。

\section{UEBA对攻击方生态识别和逼移感知的增益}

\subsection{攻击方生态识别}

UEBA行为基线为区分攻击方类型提供了更精细的依据。
黑产级攻击方使用自动化工具，行为模式单一，
与正常用户基线有显著偏离——偏离幅度大但持续时间短。
战略级攻击方使用定制工具和被盗凭据，
行为模式更接近正常用户基线——偏离幅度小但长期存在。

UEBA通过偏离度幅度和持续时间两个维度，自动区分两类攻击方。
识别结果触发自适应策略切换——
黑产级触发白帽子瓦解策略，战略级触发国家级防御响应。

\subsection{逼移时机感知}

原有的逼移时机感知依靠响应延迟和封禁策略相变来推断蓝队饱和状态。
UEBA提供了一种更灵敏的感知手段。
攻击方从广泛扫描转变为集中逼移时，
其行为基线从扫描基线急剧转变为逼移基线。
UEBA不仅知道“探测强度在上升”，
更知道这种上升是否伴随行为模式的根本性转变——
而行为模式的根本性转变，是逼移送信号即将来临的前兆。

\subsection{蜜罐引流时机裁决}

当UEBA检测到攻击方行为模式从扫描基线急剧转变为漏洞利用基线时，
说明攻击方即将执行关键操作。
此时引流时机恰好是行为基线发生根本性转变的时刻：
防御方应在攻击方完成转变之前，将其引向蜜罐，
最大化TTP收集收益。

\section{UEBA基线参数的柯克霍夫保护}

UEBA基线的具体参数——基线偏离阈值、正常行为基线的维度权重、
异常评分的聚合方式——需要保密。
这些参数本质上告诉攻击方“什么样的行为会被视为异常”。
如果攻击方掌握了这些参数，可以校准自己的行为以规避检测。

UEBA基线参数需要周期性扰动。
每个周期微调基线偏离阈值，改变正常行为基线的某些维度权重。
攻击方在上一个周期校准的行为模式，在新周期可能触发告警。
这让攻击方无法通过长期潜伏来完全规避行为检测。

\textbf{单周期内自适应校准的对抗策略}：
周期性扰动对抗的是攻击方跨周期积累知识的探测。
然而，攻击方可能在同一周期内通过逐步试探来逼近阈值边界——
每次试探后立即调整行为，
不依赖跨周期知识，只依赖当前周期内“做了什么还没被检测”的实时反馈。
周期性扰动对此类攻击无效。

对抗单周期内自适应校准的策略，是在基线偏离阈值上叠加随机噪声。
每次对同一行为的判定结果带有不可消除的随机性——
攻击方即使做了某件事没触发告警，
也不能确定下一次做同样的事是否仍然安全。
周期性扰动对抗跨周期校准，随机噪声对抗单周期内试探。
两者叠加构成对抗自适应探测的完整防线。

UEBA基线参数的保密、周期性扰动和随机噪声策略，
与已有的参数保密框架完全同构——
算法公开，行为基线参数保密，周期性轮换，实时噪声。
这是柯克霍夫原则在行为分析领域的完整应用。

\section{UEBA误报的负反馈循环与最优筛选阈值}

UEBA越灵敏，产生的基线偏离事件越多。
这些事件中相当一部分是误报——
正常用户的合法行为波动被误判为异常。
如果分析能力被UEBA自身的误报淹没，
系统将丧失检测真正攻击的能力。
这是一个危险的负反馈循环：
引入UEBA本意是增强对手行为分析，
但其误报反而削弱了告警分析层面。

为避免这一循环，UEBA产生的基线偏离事件不应全部触发人工研判。
大部分低偏离度事件直接归档，
仅在高偏离度事件上触发人工分析。
偏离度筛选阈值的最优设定，
是边际检测收益等于边际误报成本的均衡点。

这正是够用就行原则在UEBA告警筛选上的应用——
不追求捕获所有基线偏离事件，
只捕获那些偏离程度足够高、边际检测收益超过误报成本的事件。
最优筛选阈值由贪心择优自动裁定。

\section{防护层自噬效应——一个普遍规律}

UEBA误报消耗分析能力进而削弱防守的现象，
并非孤立存在于行为分析模块中。
它在安全参数自适应调节中出现过——
战时加密审计消耗算力挤占告警分析。
在零知识合规审计中出现过——
审计证明生成成本挤占告警分析。
在蜜罐运维中出现过——
蜜罐越真实，运维消耗越多，可用于告警分析的能力越少。

这些循环的共同结构是一个普遍的安全困境——
当一个防护层消耗的分析能力超过了它所能防止的损失时，
它就变成了系统的负担。

够用就行原则是克制自噬效应的唯一手段。
最优防护强度恰好是防护层的边际防护收益
等于其边际分析能力消耗的均衡点。
超过此点，防护层仍在保护系统，
但它消耗的能力已超过它所能防止的损失，
系统在净收益上为负。
这个规律值得被提炼为独立的方法论概念——防护层的自噬效应。

\section{诚实标注}

\begin{center}
\begin{xltabular}{\textwidth}{c|X|c}
\toprule
\textbf{命题} & \textbf{状态} & \textbf{层级} \\
\midrule
UEBA基线粒度的成本函数 &
二次成本函数假设合理性需实测检验，
U型安全损失曲线的成立依赖漏报与误报损失的可比量级 &
II（定性正确，精确参数待标定） \\
最优基线粒度的边际均衡条件 &
形式化推导已完成，均衡条件明确 &
II \\
U型安全损失函数的边界条件 &
当漏报损失远大于误报损失时，最优基线粒度趋向极细端，
此时计算和存储成本成为限制因素 &
II（已识别，定量分析待后续展开） \\
UEBA运维成本的多维拆分 &
单一系数在原型阶段适用，
工程实现时应拆分为算力、存储、带宽三个子项 &
II（技术注释，工程实施时独立标定） \\
UEBA对攻击方生态识别的增益 &
偏离度-持续时间双维度区分逻辑清晰，
实际区分精度需标注样本验证 &
III \\
UEBA对逼移时机感知的增益 &
行为基线转变作为逼移送信号前兆的逻辑成立 &
III（实战效果待验证） \\
UEBA对蜜罐引流时机的增益 &
行为基线转变作为引流触发信号的逻辑成立 &
III（实战效果待验证） \\
UEBA基线参数的柯克霍夫保护 &
周期性扰动对抗跨周期探测，随机噪声对抗单周期试探，
两者叠加构成完整防线 &
II（双策略逻辑完备） \\
UEBA误报的负反馈循环与最优筛选阈值 &
已识别，贪心择优处理逻辑已明确 &
II（最优筛选阈值公式待推导） \\
防护层自噬效应 &
识别为多个安全能力共享的普遍规律，
够用就行原则是克制自噬的唯一手段 &
II（定性规律已确立，形式化表达待展开） \\
UEBA与四层分配框架的对接 &
耦合关系已明确，最优分配由贪心择优统一裁定 &
II \\
\bottomrule
\end{xltabular}
\end{center}

\section{结语}

UEBA的引入将对手行为分析从统计信号的抽象层
降落到用户和实体的具体行为基线。
最优基线粒度是够用就行原则在行为分析维度的直接表达——
既要灵敏到捕捉攻击方的行为转变，
又不能灵敏到被自身的误报淹没。
UEBA基线参数的柯克霍夫保护
通过周期性扰动和随机噪声的双重机制，
对抗攻击方的跨周期探测和单周期试探。
防护层自噬效应的识别揭示了
引入任何安全能力都必须回答的元问题——
它保护了多少，又消耗了多少。
够用就行原则是唯一的度量衡。

\vspace{5mm}
\noindent\textbf{文档信息}：
\begin{itemize}
    \item 作者：李龙胜
    \item 版本：修正版（整合外部审查反馈）
    \item 许可：CC BY 4.0
    \item 前置依赖：四卷体系、安全参数自适应调节文档、
          蜜罐运维与反制模块、产品蓝图修正版
\end{itemize}

\end{document}